\section{The single upstream closure object: the physical $n=2$ mixed Hessian}
\label{sec:bhsm-n2-mixed-hessian}

The finite-kernel inventory shows that the five anisotropic observable kernels
are not presently independent numerical objects.  Their common upstream
input is the physical metric--topographic Hessian in the first admissible
$n=2$ scalar sector.  This section reduces that object as far as the retained
BHSM action presently permits.

\subsection{Retained scalar/topographic operator}

The BHSM v5.11 quadratic-operator audit gives the formal scalar/topographic
block
\begin{equation}
\mathcal H_{\rm ST}
=
\begin{pmatrix}
-\partial_\rho^2+5 & -\mathcal I_{\rm mix}\\
-\mathcal I_{\rm mix}^\dagger & -\Delta_B+5
\end{pmatrix},
\label{eq:bhsm-ST-formal}
\end{equation}
where the collar field and Berger-boundary field require a trace/extension
map $\mathcal I_{\rm mix}$ whose full nonhomogeneous domain remains open.

For the physical scalar harmonic $n=2$ on a unit $S^3$ section,
\begin{equation}
-\Delta_B e_2=8e_2.
\end{equation}
Let $\lambda_\rho\ge0$ denote the collar radial eigenvalue on the selected
self-adjoint radial domain and let
\begin{equation}
I_2
=
\langle u_\rho,\mathcal I_{\rm mix}e_2\rangle
\end{equation}
be the projected trace/extension matrix element.  Then the retained
scalar/topographic $n=2$ block has the conditional form
\begin{equation}
\boxed{
\mathcal C_{2,\rm ST}
=
\frac1{L^2}
\begin{pmatrix}
\lambda_\rho+5 & -I_2\\
-I_2^\ast & 13
\end{pmatrix}.
}
\label{eq:C2-ST}
\end{equation}

\paragraph{Theorem 7 (scalar/topographic $n=2$ positivity criterion).}
On a one-mode collar/boundary reduction with
$\lambda_\rho\ge0$, the block in Eq.~\eqref{eq:C2-ST} is positive definite if
and only if
\begin{equation}
\boxed{
|I_2|^2
<
13(\lambda_\rho+5).
}
\label{eq:I2-positivity}
\end{equation}

\paragraph{Proof.}
The first principal minor is $\lambda_\rho+5>0$.  Positivity of the
$2\times2$ Hermitian matrix is therefore equivalent to positivity of its
determinant,
\[
13(\lambda_\rho+5)-|I_2|^2>0.
\]
The common factor $L^{-2}$ is positive. \hfill$\square$

\subsection{Physical metric--topographic block}

Let $\mathcal A_g^{(2)}$ denote the gauge-reduced scalar metric Hessian
projected onto the same physical $n=2$ representation, and let
$\mathcal B_{g2}$ denote the projected metric--topographic mixed block.  The
complete retained cosmological closure object is
\begin{equation}
\boxed{
\mathbb H_2^{\rm phys}
=
\begin{pmatrix}
\mathcal A_g^{(2)} & \mathcal B_{g2}\\
\mathcal B_{g2}^\dagger & \mathcal C_{2,\rm ST}
\end{pmatrix}.
}
\label{eq:H2-phys}
\end{equation}

The later BHSM action audit already classifies the metric--scalar/topographic
mixing as nonzero in source structure but not yet fully evaluated on a common
nonhomogeneous domain.  Therefore neither
$\mathcal A_g^{(2)}$ nor $\mathcal B_{g2}$ may be replaced by the positive
diagnostic coefficients used in the finite gauge-bookkeeping model.

\subsection{Single upstream closure theorem}

\paragraph{Theorem 8 (one Hessian closes all five observable kernels).}
Assume that:

\begin{enumerate}
\item $\mathbb H_2^{\rm phys}$ is evaluated from the retained action on a
      common self-adjoint physical domain;
\item $\mathcal A_g^{(2)}$ is invertible after the physical gauge quotient;
\item Eq.~\eqref{eq:I2-positivity} and the full Schur-positivity conditions
      hold on the selected branch;
\item the initial-data prescription for the reduced $n=2$ branch is fixed.
\end{enumerate}

Then the five R1 observable kernels
\begin{equation}
\mathcal F_{D_L},\quad
\mathcal F_H,\quad
\mathcal F_{D_A},\quad
\mathcal F_{f\sigma_8},\quad
\mathcal F_{\rm lens}
\end{equation}
are all determined by this single upstream object together with covariant
matter conservation and the observation functionals of
Sec.~\ref{sec:bhsm-observable-transfer}.

\paragraph{Proof.}
The physical Schur reduction gives
\begin{equation}
\mathcal K_2
=
\mathcal C_{2,\rm ST}
-
\mathcal B_{g2}^\dagger
(\mathcal A_g^{(2)})^{-1}
\mathcal B_{g2},
\end{equation}
and the induced metric response is
\begin{equation}
\mathcal R_{g2}
=
-(\mathcal A_g^{(2)})^{-1}\mathcal B_{g2}.
\end{equation}
The matter-response theorem then fixes $\mathcal R_{m2}$ from
$\mathcal R_{g2}$ and the same $n=2$ state.  Finally each observable kernel is
the corresponding linear observation functional acting on the common
response.  No independent kernel dynamics remain. \hfill$\square$

\subsection{Correction to the numerical inventory}

The existing \texttt{code/growth.py} quantity
\begin{equation}
R_{f\sigma_8}
=
D'_{\rm R1}/D'_{\Lambda{\rm CDM}}
\end{equation}
is a useful isotropic/reference growth-ratio diagnostic.  It is not the
directional BHSM-native kernel $\mathcal F_{f\sigma_8}$ defined in the present
closure protocol.  It must therefore not be promoted to that role without
the common $n=2$ response map.

Likewise, the current luminosity-distance, Hubble and lensing expressions in
the manuscript specify observation functionals or phenomenological
relationships, not action-evaluated anisotropic kernels.

\subsection{Scientific consequence}

The numerical problem is therefore smaller than ``compute five unrelated
kernels'' but more fundamental than the existing background-growth
calculation.  The exact remaining action object is
\begin{equation}
\boxed{
\left\{
\mathcal A_g^{(2)},
\mathcal B_{g2},
I_2,
\lambda_\rho,
\mathcal D_{\rm phys}^{(2)}
\right\},
}
\label{eq:remaining-Hessian-data}
\end{equation}
where $\mathcal D_{\rm phys}^{(2)}$ is the common self-adjoint physical
domain.

Once Eq.~\eqref{eq:remaining-Hessian-data} is action-derived, all five
observable kernels follow through Theorems 1--8 without introducing new
anisotropic fit parameters.
